\documentclass[11pt]{article}
\usepackage{amsmath,amssymb,amsthm,hyperref}
\newtheorem{proposition}{Proposition}
\begin{document}
\title{A polynomial-syzygy consequence of the existing AC2 upper bound}
\author{}\date{}\maketitle
This note gives a structural consequence of the established spin-two
inertia argument. It does not improve the unrestricted upper bound.

Identify the spin-$N/2$ space with homogeneous binary forms of degree $N$,
using the normalized Dicke map
\[
 (c_0,\ldots,c_N)\longmapsto
 f(x,y)=\sum_{r=0}^N c_r\sqrt{\binom Nr}\,x^{N-r}y^r.
\]
Let $\mathcal C$ be a $k$-dimensional second-order AC subspace, $N\ge4$.
\begin{proposition}
The multiplication map
\[
 \mathcal C\otimes\operatorname{Sym}^4\mathbb C^2
 \longrightarrow\operatorname{Sym}^{N+4}\mathbb C^2,
 \qquad f\otimes g\longmapsto fg,
\]
is injective. In particular, if the forms in $\mathcal C$ have a common
factor of degree $g$, then
\[
 \boxed{5k\le N+5-g.}
\]
\end{proposition}
\begin{proof}
Put $j=N/2$, let $S_a$ be auxiliary spin-two matrices, and set
\[
 X=\sum_aJ_a\otimes S_a,\qquad
 W=X^2+(j+3)X-2j(j+1)I.
\]
The AC2 equations imply that the compression of $W$ to
$\mathcal C\otimes V_2$ is zero. In the total-spin sectors $j+r$,
$r=-2,-1,0,1,2$, the eigenvalues of $W$ are respectively
\[
 -2(j+1),\quad -2j(j+1),\quad -j(2j+5),\quad -5j-2,
 \quad4j(j+1).
\]
Thus $W$ is strictly negative on the kernel of the highest-spin
projection. A nonzero vector of $\mathcal C\otimes V_2$ cannot belong
to that kernel. The highest-spin projection, identified with the Cartan
product, is the displayed multiplication map up to nonzero normalization.
This proves injectivity.

If $h$ is a common factor of degree $g$, every product belongs to
$h\operatorname{Sym}^{N+4-g}\mathbb C^2$, of dimension $N+5-g$.
The domain has dimension $5k$, proving the bound.
\end{proof}
Equivalently, for a basis $f_1,\ldots,f_k$ of $\mathcal C$, there is no
nontrivial polynomial syzygy
\[
 \sum_{a=1}^k f_a g_a=0
\]
whose coefficients $g_a$ are homogeneous forms of common degree at most
four. A relation of smaller degree can be multiplied by a nonzero form
to obtain a degree-four relation.

\paragraph{A fixed-parity corollary.}
If $N$ is odd and every code polynomial has only even powers of $y$,
all forms are divisible by $x$. If every form has only odd powers of
$y$, all are divisible by $y$. Thus any code contained in a single
Dicke-parity sector satisfies
\[
 5k\le N+4.
\]
For $N=10m+5\ge35$ this lowers the current general integer upper bound
by one within either such sector: $k\le2m+1$, rather than $2m+2$.
It does not apply to arbitrary codes, or to a code merely invariant under
parity while using both sectors.

For arbitrary subspaces the common factor may have degree zero, so this
recovers only the current unrestricted bound $5k\le N+5$.
Injectivity with degree-five coefficients would imply $6k\le N+6$,
but no such result is proved here. A degree-five relation can involve at
most six independent code vectors after changing the code basis, since
$\dim\operatorname{Sym}^5\mathbb C^2=6$. Therefore that proposed extension
reduces to a pure-state problem in $V_j\otimes V_{5/2}$ with an AC2
physical Schmidt support of dimension at most six. A mixed-state
relaxation of this problem is insufficient: the lower total-spin sectors
can support mixtures having the required scalar first and second moments.
This is a limitation of that relaxation, not a proof that a pure solution
exists or is impossible.
\end{document}
